\documentclass[10pt,a4paper]{amsart}
\usepackage[margin=19mm]{geometry}
\usepackage{amsmath,amssymb,mathtools}
\usepackage[scr=rsfs,cal=euler]{mathalfa}
\usepackage{xfrac}
\usepackage[unicode,hidelinks,bookmarks=false]{hyperref}
\newcommand{\ie}{i.e.\ }
\newcommand{\cf}{cf.\ }
\newcommand{\resp}{respectively\ }
\newcommand{\Subsection}{\subsection}
\providecommand{\nobreakdash}{\nobreak\hbox{-}}
\setlength{\emergencystretch}{2em}
\multlinegap0pt
\usepackage{bm}
\usepackage{bbm}
\usepackage{dsfont}
\usepackage{xspace}
\usepackage{cancel}
\usepackage{mathtools}
\usepackage{amsthm}
\makeatletter
\@ifundefined{theorem}{\newtheorem{theorem}{Theorem}}{}
\@ifundefined{lemma}{\newtheorem{lemma}[theorem]{Lemma}}{}
\makeatother
\title[Spectral Radius, Vertex Deletion, and Chromatic Number of Si]{Spectral Radius, Vertex Deletion, and Chromatic Number of Signed Graphs}
\author{Abhay Jayarajan, M. Rajesh Kannan, Priti Prasanna Mondal, Shivaramakrishna Pragada}
\date{Source: arXiv:2606.23584v3 (CC BY 4.0)}
\begin{document}
\maketitle

\noindent\emph{Source and licence.} Spectral Radius, Vertex Deletion, and Chromatic Number of Signed Graphs. Version arXiv:2606.23584v3 (CC BY 4.0), distributed under the Creative Commons Attribution 4.0 International licence, as recorded in the arXiv licence metadata for this exact version and shown on the abstract page https://arxiv.org/abs/2606.23584v3. Full rights notice: licenses/arxiv-2606.23584-CC-BY-4.0.txt.\par\smallskip

\noindent\textit{Editorial scope.} Counted unit: the Theorem whose source label is \texttt{Spec\_Rad\_Bound1} in the article (source0.tex lines 592--600, source label \texttt{Spec\_Rad\_Bound1}), delivered together with the article's own complete original proof of it (source0.tex lines 601--726, one \texttt{proof} environment of 126 lines). No mathematics is written, repaired or re-derived: statement and proof are the article's own text line for line. Required context retained uncounted: Spec\_Rad\_Theta (lines 505--515); Spec\_Rad\_lower (lines 566--570). Every definition, display and label of those lines is retained unchanged. Omitted from this package and not redistributed: 1--504, 516--565, 571--591, 727--1772. Nothing inside a retained excerpt is omitted or altered: every retained body is delivered exactly as the article's lines are written, so no omission receipt of any kind accompanies this package. Citations: the retained excerpts carry no citation command at all, so no bibliography entry of the article is transcribed and no reference is redistributed. Reference adaptation: none was needed and none was made; every reference inside the delivered unit resolves inside it, so no unresolved reference or citation is produced. Printed slips kept literal: no printed word, symbol, hypothesis or number of the counted statement or of its proof was altered. Layout and class: the article's own class is \texttt{article} with packages including fullpage, amsmath, amsthm, amsfonts, amssymb, amstext, mathrsfs, enumerate, mathtools, graphicx, ragged2e, lscape, framed, xcolor; the wrapper is the lane's pinned \texttt{amsart} editorial wrapper at 10pt on A4, which restates the theorem-like environments the retained text needs and adds the article's own macro definitions that the retained text needs (none), with extra wrapper packages (bm, bbm, dsfont, xspace, cancel, mathtools). The delivered numbering is the unit's own. Assets: no retained span contains a figure environment, a graphics-inclusion command, a PDF-page inclusion, a table or a TikZ picture; no image file is copied into this package, the package declares no asset dependency and no image file is redistributed with it. Licence: version arXiv:2606.23584v3 of this article is distributed under the Creative Commons Attribution 4.0 International licence, as recorded in the arXiv licence metadata for this exact version and shown on the abstract page https://arxiv.org/abs/2606.23584v3. A full rights notice is delivered at licenses/arxiv-2606.23584-CC-BY-4.0.txt. Characters: every retained excerpt is pure ASCII, so the delivered engine, XeLaTeX with its default Latin Modern text font, reports no missing character.
	\begin{theorem}\label{Spec_Rad_Theta}
		Let $\Sigma=(G,\sigma)$ be a signed graph and let $v$ be a non-isolated vertex. Let $\theta^{*}$ denote the largest real root of $g_v(x)$.
		Then
		\[
		\rho^{2}(\Sigma)\leq \rho^{2}(\Sigma-v)+\theta^{*}.
		\]
		Furthermore, if $t_v^+=t_v^-$, then
		\[
		\rho^{2}(\Sigma)\leq \rho^{2}(\Sigma-v)+d(v).
		\]
	\end{theorem}

	\begin{lemma}\label{Spec_Rad_lower}
		For $n \geq 2$,
		\[\rho^2(\Sigma)\geq \dfrac{4(\vert E^+(\Sigma)\vert-\vert E^-(\Sigma)\vert)^2}{(n-1)^2}-(2n-1).\]
		Equality holds if and only if $\Sigma\cong(K_n,+)$ or $\Sigma \cong(K_n,-)$.
	\end{lemma}

	\begin{theorem}\label{Spec_Rad_Bound1}
		Let $\Sigma=(G,\sigma) $ be a signed graph and $v$ be a non-isolated vertex.  Then \[\rho(\Sigma) \leq \sqrt{\rho^2(\Sigma-v)+2d(v)-1}.\]
		Let $\Sigma^\prime$ be the signed graph obtained from  $\Sigma$ by switching with respect to $N^{-}(v)$. Let $d=d(v) \ge 2$, then the equality holds if and only if
		\begin{itemize}
			\item[(a)] $\Sigma^\prime[N(v)]\cong (K_d,+)$ or
			$\Sigma^\prime[N(v)]\cong (K_d,-)$; and
			\item[(b)] $\rho(\Sigma-v)=d-1$.
		\end{itemize}
	\end{theorem}

	\begin{proof}
		Let $\Sigma^\prime$ be the signed graph obtained from $\Sigma$ by switching with respect to the set $N^-(v)$. Then the adjacency matrix of $\Sigma^\prime$ can be written in the block form
		\[
		A(\Sigma^\prime)=
		\begin{pmatrix}
			0 & a^T\\
			a & B
		\end{pmatrix},
		\]
		where $B=A(\Sigma^\prime-v)$ and $a$ is the column vector whose entries indicate the neighbours of $v$ in $\Sigma^\prime-v$.
		
		Let $\theta^*$ be the largest real root of $g_v(x)$.
		If $d(v)=1$, then $t_v^+=t_v^-=0$, and hence we must have
		\[
		g_v(x)=(x-d(v))^2\bigl(\rho^{2}(\Sigma-v)+x\bigr).
		\]
		Since $g_v(d(v))=0$, $\theta^*=d(v)$, and the result follows immediately.
		
		So assume that $d(v)\ge 2$. Since $2d(v)-1>d(v)$, Theorem~\ref{Spec_Rad_Theta} implies that it suffices to show that $g_v(2d(v)-1)\ge 0.$
		We have
		\begin{align}
			\rho^2(\Sigma-v)
			&=\rho^2(\Sigma^\prime-v) \notag\\
			&\ge \rho^2(\Sigma^\prime[N(v)]) \label{Nbr_Ineq_1} \\
			&\ge
			\frac{4\bigl(|E^{+}(\Sigma^\prime[N(v)])|-|E^{-}(\Sigma^\prime[N(v)])|\bigr)^2}
			{(d(v)-1)^2}
			-(2d(v)-1) \label{Nbr_Ineq_2}\\
			&=
			\frac{4\vert q_v(\Sigma)\vert^2}{(d(v)-1)^2}
			-(2d(v)-1)\notag,
		\end{align}
		where the second inequality follows from Lemma~\ref{Spec_Rad_lower}. The last equality follows from the fact that 
		\[
		|E^{+}(\Sigma^\prime[N(v)])|-|E^{-}(\Sigma^\prime[N(v)])|
		= q_v(\Sigma')=q_v(\Sigma).
		\]
		
		Multiplying the above inequality by $(d(v)-1)^2$ yields
		\[
		(d(v)-1)^2\bigl(\rho^2(\Sigma-v)+2d(v)-1\bigr)
		\ge 4\vert q_v(\Sigma)\vert^2.
		\]
		Since
		\[
		g_v(2d(v)-1)
		=(d(v)-1)^2\bigl(\rho^2(\Sigma-v)+2d(v)-1\bigr)
		-4\vert q_v(\Sigma) \vert^2,
		\]
		and $d(v)-1\ge 1$, it follows that
		\[g_v(2d(v)-1)\ge 0.\]
		Therefore, $\theta^*\le 2d(v)-1$. Combining this with Theorem~\ref{Spec_Rad_Theta} completes the proof of the inequality. 
		
		For the equality, let $d(v) \geq 2$. Suppose first that
		\[ \rho^2(\Sigma)=\rho^2(\Sigma-v)+2d(v)-1. \]
		From the proof of the inequality,
		\[ \rho^2(\Sigma) \leq \rho^2(\Sigma-v)+\theta^* \leq \rho^2(\Sigma-v)+2d(v)-1. \]
		Hence $\theta^*=2d(v)-1$, and therefore
		\[ g_v(2d(v)-1)=0. \]
		It follows that
		\[ \rho^2(\Sigma-v) = \frac{4\vert q_v(\Sigma)\vert^2}{(d(v)-1)^2} -(2d(v)-1).\]
		Together with \eqref{Nbr_Ineq_1} and \eqref{Nbr_Ineq_2}, this gives
		\[ \rho^2(\Sigma-v) =\rho^2(\Sigma^\prime-v) \geq \rho^2(\Sigma^\prime[N(v)])
		\geq \frac{4\vert q_v(\Sigma)\vert^2}{(d(v)-1)^2} -(2d(v)-1) =\rho^2(\Sigma-v).\]
		Thus equality holds in both \eqref{Nbr_Ineq_1} and \eqref{Nbr_Ineq_2}. By Lemma~\ref{Spec_Rad_lower}, condition (a) follows. Equality in \eqref{Nbr_Ineq_1} then gives
		\[ \rho^2(\Sigma-v) =\rho^2(\Sigma^\prime[N(v)]) =(d(v)-1)^2, \]
		and hence condition (b) follows.
		
		Conversely, suppose that conditions (a) and (b) hold. By condition (b) and switching invariance,
		\[ \rho(B)=\rho(\Sigma^\prime-v)=\rho(\Sigma-v)=d(v)-1. \]
		
		First suppose that $\Sigma^\prime[N(v)]\cong(K_d,+)$. Then
		\[ a^TBa=d(v)(d(v)-1) \quad\text{and}\quad a^Ta=d(v). \]
		
		Thus the Rayleigh quotient of $a$ for $B$ is $d(v)-1=\rho(B)$. Since $B$ is symmetric, equality in the Rayleigh quotient gives
		\[ Ba=(d(v)-1)a. \]
		Consequently,
		\[ 
		A(\Sigma^\prime)
		\begin{pmatrix}
			1\\
			a
		\end{pmatrix}
		=
		\begin{pmatrix}
			a^Ta\\
			a+Ba
		\end{pmatrix}
		=d(v)
		\begin{pmatrix}
			1\\
			a
		\end{pmatrix}. \]
		Hence $d(v)$ is an eigenvalue of $A(\Sigma^\prime)$.
		
		Now suppose that $\Sigma^\prime[N(v)]\cong(K_d,-)$. Then
		\[ a^TBa=-d(v)(d(v)-1) \quad\text{and}\quad a^Ta=d(v). \]
		Thus the Rayleigh quotient of $a$ for $B$ is $-(d(v)-1)$, whose absolute value is $\rho(B)$. Since $B$ is symmetric, equality in the Rayleigh quotient gives
		\[ Ba=-(d(v)-1)a. \]
		Consequently,
		\[
		A(\Sigma^\prime)
		\begin{pmatrix}
			-1\\
			a
		\end{pmatrix}
		=
		\begin{pmatrix}
			a^Ta\\
			-a+Ba
		\end{pmatrix}
		=-d(v)
		\begin{pmatrix}
			-1\\
			a
		\end{pmatrix}.\]
		Hence $-d(v)$ is an eigenvalue of $A(\Sigma^\prime)$.
		
		Therefore, in either case,
		\[ \rho(\Sigma)=\rho(\Sigma^\prime)\geq d(v). \]
		On the other hand, the inequality already proved and condition (b) give
		\[ \rho(\Sigma) \leq \sqrt{(d(v)-1)^2+2d(v)-1} =d(v).\]
		Thus
		\[\rho^2(\Sigma)=\rho^2(\Sigma-v)+2d(v)-1,\]
		which completes the equality case.
	\end{proof}

\end{document}
